%% kapa: exam=1 round=23 year=2012 period=1 form=A subject=economics subject_ko=경제원론 seq=15 number=55
%% source: Q-Net 감정평가사 기출문제 (cst00309) artlSeq=5052578 file=raw/1008/5052578/2012년 제23회 감정평가사 제1차시험 기출문제(1교시 A형).pdf
%% conversion: kordoc 4.24.51 -> tools/gen1.py (auto); review=proofread-sonnet; flags=none
\begin{question}{55}
甲과 乙 두 사람이 사는 사회에서 甲의 소득을 X, 乙의 소득을 Y라 표시하고, 이들의 소득 분포는 (X, Y)의 형태로 표시한다. 소득 분포 상태를 평가하는 세 가지 원칙은 아래와 같다. 다음 설명으로 옳지 않은 것은?
\begin{framedtext}
○ A : 사회에서 가장 가난한 사람의 소득이 높을수록 바람직하다.\\
○ B : 모든 사회 구성원들의 소득의 총합이 클수록 바람직하다.\\
○ C : 모든 사회 구성원들의 소득이 균등하게 분포될수록 바람직하다.
\end{framedtext}
\begin{choices}
  \item 소득분포 (3, 2)와 (5, 1)을 비교할 때, 원칙 A에 따르면 (3, 2)가 더 바람직하다.
  \item 소득분포 (3, 2)와 (4, 2)를 비교할 때, 원칙 B에 따르면 (4, 2)가 더 바람직하다.
  \item 소득분포 (1, 1)과 (4, 1)을 비교할 때, 원칙 C에 따르면 (1, 1)이 더 바람직하다.
  \item 소득분포 (3, 3)과 (2, 3)을 비교할 때, 위 세 가지 원칙 모두 (3, 3)을 더 바람직하다고 판단한다.
  \item 소득분포 (2, 3)과 (7, 3)을 비교할 때, 위 세 가지 원칙 중 (7, 3)이 명백히 더 바람직하다고 판단하는 원칙은 B 뿐이다.
\end{choices}
\end{question}
\answer{5}
